\begin{textsample}{NEG Dim 3 – human – Score 1.00 – es/es\_ef\_anos\_finais\_linguagens\_ingles\_2\_p19.txt}  \label{ex:f3_neg_007}
 TEMAS INTEGRADORES : compreendem aspectos que ultrapassam a dimensão cognitiva, de modo a promover a compreensão dos processos sociais e a mudança de comportamentos que comprometem a convivência democrática. É facultado ao professor integrar aos seus planos de aula outros temas que julgue ser convenientes, considerando as especificidades de seus contextos.
Outro aspecto importante do Currículo do Espírito Santo que deve ser evidenciado no planejamento das aulas de Língua Inglesa é a intencionalidade pedagógica nos processos de desenvolvimento das competências socioemocionais, o que potencializa a construção do repertório linguístico discursivo do estudante, promovendo sua autonomia, a elevação de sua autoestima e autoconfiança, na medida em que oferece possibilidades de aprender a conhecer, aprender a ser, aprender a fazer e aprender a conviver por meio de abordagens que o levem a desenvolver o pensamento crítico, autônomo, participativo, solidário e atuante em todo o processo educacional.

% matched lemmas: 
\end{textsample}
